% Part: lambda-calculus
% Chapter: introduction
% Section: syntax

\documentclass[../../../include/open-logic-section]{subfiles}

\begin{document}

\olfileid[tr]{lam}{int}{syn}
\olsection{Lambda Kalkülüsünün Sözdizimi}

$x$, $y$, $z$,~\dots değişken dizisiyle ve $a$, $b$, $c$,~\dots sabit
simgeleriyle başlarız. Terimler kümesi tümevarımlı olarak şöyle tanımlanır:
\begin{enumerate}
\item Her değişken bir terimdir.
\item Her sabit bir terimdir.
\item $M$ ve $N$ terimse $(MN)$ de terimdir.
\item $M$ bir terim ve $x$ bir değişkense $(\lambd[x][M])$ de terimdir.
\end{enumerate}
$(MN)$ biçimindeki terimlere \emph{uygulamalar}, $(\lambd[x][M])$
biçimindekilere \emph{soyutlamalar} denir.

Hiç sabit içermeyen sisteme \emph{salt} lambda kalkülüsü denir. Çoğunlukla
salt $\lambd$-kalkülüsünde çalışacağımızdan bütün küçük harfler değişkenleri
gösterecektir. $\lambd$-kalkülüsünün terimlerini göstermek için büyük harfler
($M$, $N$ vb.) kullanırız.

Birkaç gösterim uzlaşımına uyacağız:

\begin{conv}
\begin{enumerate}
\item Parantezler yazılmadığında uygulama soldan sağa gerçekleşir. Örneğin
  $M$, $N$, $P$ ve $Q$ terimse $MNPQ$, $(((MN)P)Q)$'nun kısaltmasıdır.
\item Yine parantezler yazılmadığında lambda soyutlamasına mümkün olan en
  geniş kapsam verilir. Örneğin $\lambd[x][MNP]$,
  $(\lambd[x][((MN)P)])$ olarak okunur.
\item Bir lambda birden çok değişkeni soyutlamak için kullanılabilir.
  Örneğin $\lambd[xyz][M]$,
  $\lambd[x][\lambd[y][\lambd[z][M]]]$'nin kısaltmasıdır.
\end{enumerate}
\end{conv}

Örneğin
\[
\lambd[xy][xxyx \lambd[z][xz]]
\]
şunun kısaltmasıdır:
\[
\lambd[x][\lambd[y][((((xx)y)x)(\lambd[z][(xz)]))]].
\]
Bu uzlaşımları ezberlemelisiniz. Başta sizi çıldırtacaklardır; ama zamanla
alışacak ve bir süre sonra parantez yığınıyla uğraşmaktan daha az
çıldırtıcı olduklarını göreceksiniz.

Yalnızca bağlı değişkenlerin adları bakımından farklı olan iki terime
$\alpha$-eşdeğer denir; örneğin $\lambd[x][x]$ ile $\lambd[y][y]$.
Bunları ``aynı'' terim saymak elverişlidir; başka bir deyişle $M$ ile $N$
aynıdır dediğimizde ``bağlı değişkenlerin yeniden adlandırılmasına kadar''
aynı olduklarını da kastederiz. Bir $\lambd$'nın kapsamında bulunan
değişkenlere ``bağlı'', diğerlerine ``serbest'' denir. Önceki örnekte serbest
değişken yoktur; fakat
\[
(\lambd[z][yz])x
\]
ifadesinde $y$ ile $x$ serbest, $z$ bağlıdır.

\end{document}
